% 图 74.5 —— 由 `checks/emit_hand.py` 自动拼出（probe 精测 + prims2 图元分解）
%%
% ⚠ 这是**稿子不是成品**：跑 `checks/checkhand.py 74.5` 看指标 + 看叠合图。
%%
%% ⚠ 待核：
%   · 本图 probe 报出阴影族 1 个 —— **本脚本不生成阴影**（要照原书手写）；prims2 的 strokes 里可能含阴影线，核对时留意别画重
%   · 可疑字形 1 项（空心圈 ⊙ 常被认成字母 o）—— 见文件末
%%
\documentclass[12pt]{standalone}
\usepackage{fontspec}
\setsansfont{Arial}
\newfontfamily\mfallback{DejaVu Sans}
\usepackage{tikz}
\usetikzlibrary{arrows.meta}
\begin{document}
\begin{tikzpicture}[x=1pt, y=-1pt, line width=4.0pt, line cap=round, line join=round]

  %% ════ 填充区（prims2 的 fills）════

  %% ════ 直线段（prims2 的 strokes）════
  % （略）线段 (702.0,61.0)-(697.0,58.0) 线宽6.2 —— 是箭头头那一坨，已由箭头头代替
  \draw[line width=3.0pt] (607.0,184.8) -- (609.0,187.0);
  % （略）线段 (47.0,148.0)-(52.0,152.0) 线宽6.5 —— 是箭头头那一坨，已由箭头头代替
  \draw[line width=6.5pt] (189.0,292.0) -- (193.0,297.0);
  \draw[line width=7.0pt] (130.0,42.0) -- (133.0,36.0);
  \draw[line width=5.3pt] (770.0,293.0) -- (769.0,298.0);
  \draw[line width=6.3pt] (844.0,136.0) -- (850.0,137.0);
  \draw[line width=7.4pt] (298.0,150.0) -- (304.0,146.0);
  \draw[line width=4.0pt] (450.0,154.0) -- (430.0,156.0);
  \draw[line width=6.5pt] (764.0,287.0) -- (769.0,291.0);
  % （略）线段 (60.0,149.0)-(54.0,153.0) 线宽7.4 —— 是箭头头那一坨，已由箭头头代替
  \draw[line width=5.0pt] (305.0,148.0) -- (304.6,297.4);
  \draw[line width=5.0pt] (304.6,297.4) -- (194.0,298.0);
  \draw[line width=5.0pt] (132.0,44.0) -- (303.1,43.1);
  \draw[line width=2.9pt] (303.1,43.1) -- (305.0,45.4);
  \draw[line width=4.0pt] (305.0,45.4) -- (305.0,145.0);
  % （略）线段 (695.0,56.0)-(697.0,50.0) 线宽6.7 —— 是箭头头那一坨，已由箭头头代替
  \draw[line width=6.5pt] (842.0,137.0) -- (838.0,142.0);
  \draw[line width=5.0pt] (54.0,153.0) -- (54.7,297.7);
  \fill (49.0,153.0) -- (59.0,153.0) -- (54.1,168.0) -- cycle;
  \draw[line width=5.0pt] (54.7,297.7) -- (191.0,300.0);
  \draw[line width=4.0pt] (53.0,151.0) -- (54.0,44.0);
  \draw[line width=5.0pt] (54.0,44.0) -- (129.0,43.0);
  \draw[line width=3.0pt] (451.0,161.0) -- (431.0,162.0);

  %% ════ 曲线（prims2 的 curves —— 三段式，坐标同为 (y,x)）════
  \draw[line width=5.0pt] (695.0,57.0) .. controls (640.2,70.5) and (603.6,128.7) .. (607.0,184.8);
  \fill (690.9,54.2) -- (699.1,59.8) -- (686.5,69.4) -- cycle;
  \draw[line width=5.0pt] (609.0,189.0) .. controls (611.5,263.9) and (698.4,322.9) .. (768.0,292.0);

  %% ════ 圆 / 椭圆（probe 精测）════
  \draw (727.6,175.7) circle (123.3);

  %% ════ 实心点 ════
  \fill (802.5,82.0) circle (5.4);
  \fill (629.5,98.0) circle (5.3);
  \fill (635.5,103.5) circle (5.7);
  \fill (613.5,188.0) circle (5.0);
  \fill (651.5,266.0) circle (4.6);
  \fill (646.0,271.5) circle (5.3);

  %% ════ 标签（probe 直接给的 \node 行）════
  %% ⚠ 全部补了 fill=white：文字在最上层，否则与曲线交叉干扰
     \node[anchor=center,inner sep=0pt,fill=white,font=\fontsize{31.9}{39.8}\selectfont] at (172.6,23.6) {\textsf{\textbf{\textit{a}}}};   % box(164,13,15,18)
     \node[anchor=center,inner sep=0pt,fill=white,font=\fontsize{32.0}{40.0}\selectfont] at (723.1,31.1) {\textsf{\textbf{\textit{a}}}};   % box(714,21,16,17)
     \node[anchor=center,inner sep=0pt,fill=white,font=\fontsize{30.7}{38.3}\selectfont] at (871.8,142.8) {\textsf{\textbf{\textit{b}}}};   % box(862,130,17,23)
     \node[anchor=center,inner sep=0pt,fill=white,font=\fontsize{28.5}{35.6}\selectfont] at (1020.2,146.5) {\textsf{\textbf{2}}};   % box(1011,138,17,16)
     \node[anchor=center,inner sep=0pt,fill=white,font=\fontsize{30.0}{37.4}\selectfont] at (1000.5,155.0) {\textsf{\textbf{\textit{F}}}};   % box(991,143,19,22)
     \node[anchor=center,inner sep=0pt,fill=white,font=\fontsize{30.0}{37.5}\selectfont] at (327.8,158.2) {\textsf{\textbf{\textit{b}}}};   % box(318,146,17,22)
     \node[anchor=center,inner sep=0pt,fill=white,font=\fontsize{41.3}{51.6}\selectfont] at (961.5,162.2) {{\mfallback\textbf{−}}};   % box(941,153,21,6)
     \node[anchor=center,inner sep=0pt,fill=white,font=\fontsize{30.0}{37.5}\selectfont] at (31.8,168.2) {\textsf{\textbf{\textit{b}}}};   % box(22,156,17,22)
     \node[anchor=center,inner sep=0pt,fill=white,font=\fontsize{37.7}{47.1}\selectfont] at (960.6,169.2) {{\mfallback\textbf{−}}};   % box(941,161,21,5)
     \node[anchor=center,inner sep=0pt,fill=white,font=\fontsize{30.0}{37.5}\selectfont] at (583.8,207.2) {\textsf{\textbf{\textit{b}}}};   % box(574,195,17,22)
     \node[anchor=center,inner sep=0pt,fill=white,font=\fontsize{31.0}{38.7}\selectfont] at (753.5,329.0) {\textsf{\textbf{\textit{a}}}};   % box(745,319,15,17)
     \node[anchor=center,inner sep=0pt,fill=white,font=\fontsize{32.0}{40.0}\selectfont] at (173.1,331.1) {\textsf{\textbf{\textit{a}}}};   % box(164,321,16,17)

  % 钉死画布：否则超出的图元/控制点会把页面撑大，1:1 回比就失准
  \pgfresetboundingbox
  \path[use as bounding box] (0,0) rectangle (1040,349);
\end{tikzpicture}
\end{document}

% ⚠ 待核清单（probe 拿不准，人眼过一遍）：
%   · 未识别/跳过：⑤ 跳过/未识别的字形
